\section{Limitations, Ethics, and Roadmap}
\label{sec:limitations-ethics-roadmap}

GAWorld makes a broad set of mechanisms available in one runtime, but breadth
does not imply scientific validity or engineering completeness.  We separate
limitations of the current artifact archive, limitations of the system
implementation, and limitations inherent in representing social life through
LLM agents.  The resulting roadmap likewise separates engineering work that
improves the instrument from research work needed to justify substantive
claims.

\subsection{Evidence and reproducibility limitations}

The archived experiments are incomplete and heterogeneous.  Several contain a
single seed or one agent, treatment horizons differ, some mechanisms produced
identically zero fields, and some intended arms lack outcome files entirely.
The evidence ledger can make those conditions visible, but it cannot recover
missing treatments or provide statistical power after the fact.  The existing
cases therefore illustrate trace production, comparison mechanics, and
failure visibility rather than stable behavioral effects.  Claims about
individual consistency, macro correspondence, emergence, or counterfactual
response require purpose-built evaluation beyond the archive summarized here.

Reproducibility is also limited by model and provider dependence.  A seed does
not freeze a hosted model, its serving stack, safety filters, or routing
behavior.  Provider fallback may change the model used within a run; library,
prompt, or parser versions can change subsequent state; and concurrency can
alter call order.  Local models improve control but do not eliminate hardware,
runtime, and numerical variation.  Consequently, meaningful replication
requires more than source code and a seed: it requires resolved configuration,
prompt and sampling settings, provider and model identifiers, dependency
versions, active plugins, warnings, per-arm completion status, and raw
artifacts.  Hosted-model studies should measure replay variability rather than
promise bitwise determinism.

Runtime and monetary cost constrain the studies that can be performed.  A
persistent society may invoke generation, embedding, retrieval, reflection,
environment, and consolidation paths repeatedly across agents and ticks.
Model caching, rule-based fallbacks, configurable stages, and local providers
can reduce calls, but each changes experimental conditions.  Cost pressure can
also encourage small populations and short horizons that are inadequate for
emergence or uncertainty analysis.  Future evaluations should report call
counts, latency, failures, model routes, and estimated cost alongside simulated
scale and horizon.

\subsection{Architecture and data limitations}

The K1--K5 microkernel migration is complete for its defined scope, including
nine built-in plugins, movement validation, and audited standard interventions.
Compatibility paths nevertheless remain: memory, environment, and social
logic are not wholly plugin-owned, the Controller does not mediate every
action, and the Recorder is not the sole provenance stream.  Shared
dictionaries and the context escape hatch weaken schema enforcement and
ownership boundaries.
Third-party plugins run in process and should not be treated as sandboxed.  A
configured extension point therefore does not guarantee mechanism isolation,
security, or semantic equivalence under replacement.

The artifact layer has corresponding schema plurality.  Agent logs, long-form
state CSVs, environment timelines, memory stores, visualization traces,
economy audits, and plugin outputs use different schemas and sometimes
different notions of step or time.  This plurality preserves domain-specific
detail, but cross-subsystem analysis requires explicit joins and can silently
misread cardinality or missingness.  The discrepancy identified in the
economy case between report wording and raw CSV cardinality illustrates this
risk.  Until a versioned provenance schema and manifest bind clocks, entities,
mechanisms, and derivations, replay and analysis must remain artifact-aware.

Input profiles are another boundary.  They may be sparse, stylized,
synthetic, or assembled from data whose completeness and provenance vary.
Rich generated backstories do not repair missing observations; they can turn
uncertainty into plausible detail.  Profile coverage, source, consent,
transformation, and missingness should be documented, and downstream readers
should distinguish supplied attributes from model-generated elaboration.

\subsection{Ethical boundaries and affected communities}

Synthetic agents can create an illusion of access to perspectives that were
never actually observed.  Language models may reproduce stereotypes, suppress
within-group heterogeneity, and make marginalized identities appear more
uniform or legible than they are \citep{wang2025flatten}.  Conditioning a model
on demographic labels or a short profile does not confer lived experience,
community membership, consent, or authority.  Even when aggregate outputs
resemble survey responses in one setting, LLM agents should not be treated as
drop-in replacements for human participants \citep{dillion2023replace}.

GAWorld therefore must not substitute for engagement with people affected by
a research question or policy.  It cannot determine community priorities,
adjudicate contested values, estimate acceptable harms, or legitimate a
decision about housing, health, employment, policing, mobility, or information
governance.  For consequential applications, affected communities should help
define the question, categories, outcomes, unacceptable uses, and release
conditions; domain experts and empirical data are needed to evaluate whether
the modeled mechanisms correspond to the institutions at issue.  Simulation
may support deliberation or identify assumptions to test, but it cannot replace
participation, due process, or accountable public decision-making.

Profiles, memories, interviews, and generated social ties can also expose or
infer sensitive information.  Researchers should minimize collection, obtain
appropriate consent and permissions, separate identifiers from research
artifacts, restrict access, define retention and deletion policies, and assess
whether releasing prompts or traces could re-identify a source.  Dashboard,
relay, and environment services are local research tooling rather than
production-secure deployments; exposing them requires authentication,
encryption, authorization, audit, and abuse controls.  Generated claims about
real people should not be presented as observations, and simulated agents
should not be used to rank, target, or make decisions about identifiable
individuals.

\subsection{Engineering roadmap}

The engineering roadmap is to make current boundaries enforceable and current
artifacts easier to reconstruct.  Priorities are: (1) extend typed Controller
validation beyond movement where governance is required; (2) migrate remaining domain mechanisms incrementally behind declared
dependencies without changing their semantics silently; (3) replace shared
dictionary conventions with versioned schemas and explicit state ownership;
(4) unify trace identifiers and clock semantics through a manifest-backed
provenance layer while retaining domain-specific payloads; (5) capture model,
provider, prompt, fallback, dependency, and plugin resolution in every run;
and (6) add conformance, failure-injection, migration, security, and replay
tests.  These changes would improve composability and auditability.  They would
not, by themselves, validate an artificial society.

\subsection{Research roadmap}

The research roadmap begins only after the execution unit is reproducible.
For individual consistency, standardized repeated probes and controlled memory
perturbations should be evaluated against held-out criteria where consent and
provenance permit.  For macro correspondence, researchers should define a
sampling frame, map simulated constructs to external measures, separate
calibration from testing, and quantify subgroup error and uncertainty.  Claims
of emergence require multiple seeds, scales, topologies, and mechanism
ablations, with macro patterns traced to micro events.  Counterfactual studies
require verified intervention paths, replicated arms, pre-specified outcomes,
placebos or negative controls, and an explicit account of transportability.
Provider and prompt sensitivity should be treated as experimental factors, not
hidden implementation noise.

Finally, future studies should evaluate not only whether a simulated pattern
is reproducible, but also whether the representation and intended use are
appropriate.  Participatory problem formulation, impact assessment, and
documentation of excluded perspectives belong beside technical validation.
GAWorld's near-term role is therefore to make artificial-society hypotheses,
mechanisms, traces, and failures easier to construct and inspect---not to turn
generated populations into unearned evidence about society.
